% Defining a comprehensive LaTeX preamble
\documentclass[aspectratio=169]{beamer}
\usepackage[utf8]{inputenc}
\usepackage{xeCJK}
\usepackage{graphicx}
\usepackage{tikz}
\usepackage{amsmath}
\usepackage{array}
\usepackage{booktabs}

% Setting Chinese font
\IfFontExistsTF{Noto Serif CJK TC}{
  \setCJKmainfont{Noto Serif CJK TC}
  \setCJKsansfont{Noto Serif CJK TC}
  \setCJKmonofont{Noto Serif CJK TC}
}{\setCJKmainfont{SimSun}
  \setCJKsansfont{SimSun}
  \setCJKmonofont{SimSun}
}

% 全域字體大小調整參數 - 可以統一調整整份投影片的字體大小
\newcommand{\globalscale}{0.9}  % 調整這個數值來統一放大或縮小字體 (0.8=80%, 0.9=90%, 1.0=100%)

% 定義統一的字體大小命令
\newcommand{\contentsize}{\fontsize{9pt}{11pt}\selectfont}
\newcommand{\tablesize}{\fontsize{7.5pt}{9pt}\selectfont}
\newcommand{\smalltablesize}{\fontsize{6.5pt}{8pt}\selectfont}

% Theme setup
\usetheme{Madrid}
\usecolortheme{default}

% Custom colors
\definecolor{ameeblue}{RGB}{0,114,188}
\definecolor{ameered}{RGB}{186,24,27}
\definecolor{ameegray}{RGB}{120,120,120}
\definecolor{lightblue}{RGB}{230,240,250}
\definecolor{lightred}{RGB}{255,240,240}
\definecolor{lightgray}{RGB}{245,245,245}

% 徹底移除所有黑色背景
\setbeamercolor{structure}{fg=ameeblue}
\setbeamercolor{frametitle}{bg=ameeblue,fg=white}
\setbeamercolor{title}{bg=ameeblue,fg=white}
\setbeamercolor{titlelike}{bg=ameeblue,fg=white}

% 調色盤設定
\setbeamercolor{palette primary}{bg=ameeblue,fg=white}
\setbeamercolor{palette secondary}{bg=ameeblue!80,fg=white}
\setbeamercolor{palette tertiary}{bg=ameeblue!60,fg=white}
\setbeamercolor{palette quaternary}{bg=ameeblue!40,fg=black}

% 標題頁設定
\setbeamercolor{title page}{fg=white}
\setbeamercolor{author in head/foot}{fg=white}
\setbeamercolor{author}{fg=black}
\setbeamercolor{institute}{fg=black}
\setbeamercolor{date}{fg=black}

\setbeamercolor{title in head/foot}{fg=white}
\setbeamercolor{institute in head/foot}{fg=white}
\setbeamercolor{date in head/foot}{fg=white}

% 背景設定為白色
\setbeamercolor{background canvas}{bg=white}
\setbeamercolor{normal text}{fg=black,bg=white}

% 設定 block 的背景和邊框顏色，避免黑色背景遮擋文字
\setbeamercolor{block title}{bg=ameeblue!80,fg=white}
\setbeamercolor{block body}{bg=lightblue,fg=black}
\setbeamercolor{block title alerted}{bg=ameered!80,fg=white}
\setbeamercolor{block body alerted}{bg=lightred,fg=black}
\setbeamercolor{block title example}{bg=ameegray!60,fg=white}
\setbeamercolor{block body example}{bg=lightgray,fg=black}

% 移除 beamer 模板中的陰影效果
\setbeamertemplate{blocks}[default]
\setbeamertemplate{title page}[default][colsep=-4bp,rounded=true]
\setbeamertemplate{part page}[default][colsep=-4bp,rounded=true]

% 調整 beamer 字體大小
\setbeamerfont{normal text}{size=\fontsize{10pt}{12pt}}
\setbeamerfont{frametitle}{size=\large}
\setbeamerfont{title}{size=\Large}

\title{AMEE 2025 Conference Highlights}
\subtitle{Barcelona, Spain}
\author{謝慕揚, MD, PhD, FESC}
\institute{台大醫院新竹分院教學部副主任}
\date{2025年8月23-27日}

\begin{document}

\begin{frame}
\titlepage
\end{frame}

\begin{frame}{目錄}
\tableofcontents
\end{frame}

% 會議簡介
\section{會議簡介}

\begin{frame}{AMEE 2025概述}
\frametitle{AMEE 2025概述}

\textbf{會議基本資訊：}
\begin{itemize}
    \item \textbf{會議名稱}：An International Association for Medical Education
    \item \textbf{地點}：Barcelona, Spain
    \item \textbf{時間}：2025年8月23-27日
    \item \textbf{性質}：全球最重要的醫學教育年會之一
\end{itemize}

\vspace{0.5cm}

\textbf{會議特色：}
\begin{itemize}
    \item 來自世界各地的醫學教育工作者齊聚
    \item 分享最新教育理念、研究成果與實務經驗
    \item 涵蓋從教育理論到實務應用的各個面向
    \item 從評估方法到教師發展的完整內容
\end{itemize}
\end{frame}

% Workshop執行方式
\section{Workshop執行模式}

\begin{frame}{AMEE Workshop特色}
\frametitle{AMEE Workshop特色}

\textbf{基本架構：}
\begin{itemize}
    \item \textbf{時間長度}：每場Workshop為三個小時
    \item \textbf{分組形式}：所有參與者分組進行,每組有自己的桌次
    \item \textbf{師資配置}：每個Workshop由三位老師共同主持
    \item \textbf{演講時間}：每位老師約有15到20分鐘的演講時間
\end{itemize}

\vspace{0.5cm}

\textbf{三位老師的協作模式：}
\begin{enumerate}
    \item \textbf{共同主持}：三個老師都一起主持,共同分擔主持工作
    \item \textbf{角色分工}：
    \begin{itemize}
        \item 其中一個老師會走到學生桌子裡面去鼓勵學生參與啟發問題
        \item 台上有兩個老師會負責交替發言提供討論
    \end{itemize}
    \item \textbf{互動模式}：大家都會互相打斷,然後很客氣地跟大家一起說自己怎麼想
\end{enumerate}
\end{frame}

\begin{frame}{Workshop的獨特學習氛圍}
\frametitle{Workshop的獨特學習氛圍}

\begin{block}{Workshop的獨特學習氛圍}
\begin{itemize}
    \item 這個會議很有趣,全部都是經驗老道的老師
    \item 鼓勵其他學生（這些學生其實本身也都是經驗豐富的老師）
    \item 很多發言,都很有哲理
    \item 每隔幾分鐘就有很精彩的對話可以被聽到
\end{itemize}
\end{block}

\vspace{0.5cm}

\textbf{學習收穫：}
\begin{itemize}
    \item 體驗經驗豐富的教育者如何進行對話
    \item 觀察如何在教學中鼓勵參與和批判性思考
    \item 學習如何建立包容性的學習環境
    \item 理解協作式教學的真正意涵
\end{itemize}
\end{frame}

% 核心教育理念
\section{核心教育理念}

\begin{frame}{Constructive Alignment（建構式配合）}
\frametitle{Constructive Alignment（建構式配合）}
\contentsize

\textbf{核心概念：}由澳洲教育學者John Biggs提出,強調教學系統中各個要素之間的一致性和協調性。

\vspace{0.3cm}

\begin{table}[h]
\centering
\tablesize
\begin{tabular}{p{2.8cm}p{8.5cm}}
\toprule
\textbf{要素} & \textbf{說明} \\
\midrule
學習目標 & 明確定義學生應該學會什麼,使用具體、可測量的行為動詞,通常採用Bloom's Taxonomy分類 \\
\midrule
教學活動 & 設計能夠幫助學生達成學習目標的學習經驗,提供適當的學習機會和練習,促進學生主動參與和深度學習 \\
\midrule
評量方法 & 評量方式必須能夠測量學生是否達成學習目標,評量標準與學習目標保持一致,包括形成性評量和總結性評量 \\
\bottomrule
\end{tabular}
\end{table}

\vspace{0.3cm}

\textbf{實施原則：}
\begin{enumerate}
    \item \textbf{向後設計（Backward Design）}：先確定學習目標,再設計評量方式,最後規劃教學活動
    \item \textbf{一致性（Consistency）}：三個要素之間必須彼此呼應,避免教學內容與評量方式脫節
    \item \textbf{透明度（Transparency）}：學習目標對學生清楚明確,評量標準公開透明
\end{enumerate}
\end{frame}

\begin{frame}{Career Trajectory（職涯軌跡）}
\frametitle{Career Trajectory（職涯軌跡）}
\contentsize

\textbf{定義與特徵：}
Career Trajectory是指個人在職業生涯中所經歷的路徑、發展階段和變化過程。

\vspace{0.3cm}

\begin{table}[h]
\centering
\tablesize
\begin{tabular}{p{3.5cm}p{7.5cm}}
\toprule
\textbf{類型} & \textbf{特徵與範例} \\
\midrule
傳統線性軌跡 & 在同一領域或組織內逐步晉升,具有清晰的階層結構。例：醫師從住院醫師→主治醫師→科主任→院長 \\
\midrule
螺旋式軌跡 & 跨領域但相關的職業轉換,技能和經驗可相互轉移。例：從臨床醫師→醫學教育→醫院管理 \\
\midrule
過渡式軌跡 & 頻繁變換不同領域的工作,追求多元經驗和技能。例：醫師兼作家、顧問、創業家 \\
\midrule
專家式軌跡 & 在特定領域深度發展,成為該領域的權威專家。例：專精某種手術技術的外科醫師 \\
\bottomrule
\end{tabular}
\end{table}

\vspace{0.3cm}

\textbf{職涯發展階段：}探索期 → 建立期 → 維持期 → 轉換期 → 傳承期
\end{frame}

% 專業主義議題
\section{專業主義議題}

\begin{frame}{Weaponized Professionalism（武器化專業主義）}
\frametitle{Weaponized Professionalism（武器化專業主義）}
\contentsize

\textbf{核心概念：}將專業主義的標準和規範作為工具,用來排斥、控制或壓制某些群體,特別是邊緣化群體的現象。

\vspace{0.3cm}

\begin{table}[h]
\centering
\tablesize
\begin{tabular}{p{3.5cm}p{7.5cm}}
\toprule
\textbf{表現形式} & \textbf{具體例子} \\
\midrule
外觀和行為標準 & 對髮型、服裝的嚴格規定（如禁止某些文化髮型）,要求「標準」的溝通方式（往往反映主流文化） \\
\midrule
語言和溝通 & 將口音或方言視為「不專業」,要求特定的溝通風格,忽視文化差異對溝通方式的影響 \\
\midrule
評量和考核 & 使用帶有文化偏見的評量標準,對來自不同背景學生的雙重標準,主觀評量中的隱性偏見 \\
\bottomrule
\end{tabular}
\end{table}

\vspace{0.3cm}

\textbf{應對策略：}
\begin{itemize}
    \item \textcolor{ameered}{重新定義專業主義}：將重點放在核心價值而非表面行為
    \item \textcolor{ameered}{教育和培訓}：提高對隱性偏見的認識,培養文化勝任力
    \item \textcolor{ameered}{政策改革}：檢視現有的專業標準和評量方式
\end{itemize}
\end{frame}

\begin{frame}{Professionalism as Conversation（對話式專業主義）}
\frametitle{Professionalism as Conversation（對話式專業主義）}

\textbf{核心理念：}「Professionalism is not a prescription but should be a conversation」

\vspace{0.5cm}

\begin{columns}
\begin{column}{0.5\textwidth}
\textbf{傳統處方式專業主義}
\begin{itemize}
    \item 提供固定的行為準則和標準
    \item 強調遵循既定規範
    \item 採用「一刀切」的方式
    \item 缺乏彈性和脈絡考量
\end{itemize}
\end{column}

\begin{column}{0.5\textwidth}
\textbf{對話式專業主義}
\begin{itemize}
    \item 鼓勵開放性討論和反思
    \item 考慮不同觀點和經驗
    \item 適應不同文化和情境
    \item 促進批判性思考
\end{itemize}
\end{column}
\end{columns}

\vspace{0.5cm}

\textbf{為什麼需要對話？}
\begin{enumerate}
    \item \textbf{複雜性與脈絡性}：專業情境往往複雜且獨特,需要考量文化、社會、個人因素
    \item \textbf{多元性與包容性}：不同背景的人對專業的理解不同,避免單一文化霸權
    \item \textbf{發展性與學習性}：專業主義是持續發展的過程,透過對話促進深度學習
\end{enumerate}
\end{frame}

% 評估相關議題
\section{評估相關議題}

\begin{frame}{Contrast Effect（對比效應）}
\frametitle{Contrast Effect（對比效應）}

\textbf{定義：}在評估或判斷時,前面經歷的刺激會影響對後續刺激的感知和評價的現象。

\vspace{0.5cm}

\textbf{在Multiple Mini-Interview (MMI)情境中的表現：}

\begin{block}{MMI（Multiple Mini Interview,多站迷你面試）}
\begin{itemize}
    \item 考生需要通過多個獨立的面試站點（通常6-10個站）
    \item 每站時間較短（通常5-10分鐘）
    \item 每站由不同的面試官評分
    \item 每站評估不同的能力或情境（溝通能力、批判思考、倫理判斷等）
\end{itemize}
\end{block}

\vspace{0.3cm}

\textbf{評分者角度：}
\begin{itemize}
    \item 如果評審員剛評完一個表現優秀的考生,可能會對下一個普通表現的考生評分偏低
    \item 相反地,在評完表現較差的考生後,可能會對接下來普通表現的考生給予較高評分
\end{itemize}
\end{frame}

\begin{frame}{減少Contrast Effect的策略}
\frametitle{減少Contrast Effect的策略}

\begin{enumerate}
    \item \textbf{標準化評分標準}：使用詳細的評分量表和具體行為指標
    
    \item \textbf{評審員訓練}：讓評審員認識這種偏誤並學習如何避免
    
    \item \textbf{休息時間}：在站點間給評審員短暫的重置時間
    
    \item \textbf{多重評審}：同一站點由多位評審員評分
    
    \item \textbf{評分校準}：定期進行評審員間的一致性訓練
\end{enumerate}

\vspace{0.5cm}

\begin{alertblock}{重要提醒}
這在臨床教育評估設計中是很重要的考量因素。設計MMI時必須考慮到對比效應對評分公平性的影響。
\end{alertblock}
\end{frame}

\begin{frame}{CCC架構（臨床能力評估委員會）}
\frametitle{CCC架構（臨床能力評估委員會）}
\contentsize

\textbf{委員會組成：}
\begin{itemize}
    \item 臨床醫師代表、醫學教育專家、基礎醫學教授
    \item 住院醫師代表、外部專家（如醫學會代表）、學生代表
\end{itemize}

\vspace{0.3cm}

\textbf{主要職責：}
\begin{itemize}
    \item 評估醫學生臨床能力、制定能力評估標準
    \item 監督實習醫師表現、處理學習困難個案、決定升級或留級
\end{itemize}

\vspace{0.3cm}

\begin{table}[h]
\centering
\smalltablesize
\begin{tabular}{p{2.3cm}p{2.8cm}p{4.2cm}}
\toprule
\textbf{衝突類型} & \textbf{情況} & \textbf{緩解措施} \\
\midrule
師生關係衝突 & 委員曾是被評估學生的指導教師 & 主動揭露師生關係,完全迴避該學生評估 \\
醫院利益衝突 & 委員所屬醫院與學生實習醫院有合作或競爭關係 & 採用匿名化評估資料,引入中立醫院專家意見 \\
家族關係衝突 & 委員與學生有親屬關係 & 絕對迴避原則,完全退出評估程序 \\
經濟利益衝突 & 委員在醫學教育相關企業有投資 & 年度財務利益揭露,涉及相關產品評估時迴避 \\
\bottomrule
\end{tabular}
\end{table}
\end{frame}

% 研究方法
\section{研究方法}

\begin{frame}{Cook (2008) 研究方法框架}
\frametitle{Cook (2008) 研究方法框架}

\textbf{教育研究的三個目的：}

\begin{table}[h]
\centering
\tablesize
\begin{tabular}{p{1.8cm}p{2.3cm}p{1.8cm}p{4.2cm}}
\toprule
\textbf{目的} & \textbf{解決的問題} & \textbf{焦點} & \textbf{理由} \\
\midrule
描述 & 做了什麼？ & 活動/事件的說明 & 可重複性、透明度,讓其他人能夠重複、理解或基於既有成果進一步發展 \\
\midrule
證成 & 有效嗎？ & 評估效果 & 效能證據,提供證據證明特定過程或方法是有益的或達到預期目的 \\
\midrule
闡明 & 為什麼/如何有效？ & 機制、理論 & 指導理論、實務,深化理解、指導未來實務,並透過闡述基本原理幫助跨情境的知識轉移 \\
\bottomrule
\end{tabular}
\end{table}

\vspace{0.3cm}

\textbf{研究現況：}
\begin{itemize}
    \item \textcolor{ameered}{\textbf{證成}}研究最為常見（72\%）
    \item 其次是\textbf{描述}（16\%）
    \item \textcolor{ameered}{\textbf{闡明}}研究最少見（12\%）
\end{itemize}

\vspace{0.2cm}

\begin{alertblock}{重要提醒}
理想上應努力進行更多闡明型研究以建立更強的理論基礎。
\end{alertblock}
\end{frame}

\begin{frame}{Post-positivist與Constructivist研究典範}
\frametitle{Post-positivist與Constructivist研究典範}
\begin{table}[h]
\centering
\smalltablesize
\begin{tabular}{p{2.2cm}p{4.3cm}p{4.3cm}}
\toprule
\textbf{面向} & \textbf{Post-positivist（後實證主義）} & \textbf{Constructivist（建構主義）} \\
\midrule
本體論（現實觀） & 相信存在客觀的現實,但認為我們只能不完美地理解它。現實是「可能真實的」(probably true),而非絕對確定 & 現實是社會建構的,存在多重的、主觀建構的現實。沒有單一的「真實」現實,而是多個同等有效的現實詮釋 \\
\midrule
認識論（知識觀） & 追求客觀性,但承認完全客觀是不可能的。強調透過嚴謹的方法控制偏誤和錯誤,相信透過多重驗證可以逼近真理 & 知識是透過研究者與被研究者的互動建構出來的。價值觀念(values)不可避免地影響研究過程 \\
\midrule
方法論 & 主要使用量化方法,但也接納質性方法作為補充。重視實驗設計、隨機對照試驗,強調信度、效度、可重複性 & 主要使用質性方法（民族誌、現象學、敘事分析等）。重視脈絡、主觀經驗和意義建構 \\
\bottomrule
\end{tabular}
\end{table}

\vspace{0.3cm}

\textbf{在MMI研究中的應用：}
\begin{itemize}
    \item \textbf{Post-positivist取向}：MMI分數與客觀結果的相關性（統計分析、信度研究、預測效度驗證）
    \item \textbf{Constructivist取向}：面試官如何在社會脈絡中建構判斷（參與觀察、深度訪談、主題分析）
\end{itemize}
\end{frame}

% AI與教育工具
\section{AI與教育工具}

\begin{frame}{AI與醫學教育工具}
\frametitle{AI與醫學教育工具}

\textbf{現場示範：ChatGPT生成考試題目}
\begin{itemize}
    \item 快速產生多樣化的評估題目
    \item 根據學習目標調整難度
    \item 提供不同形式的評量方式
\end{itemize}

\vspace{0.5cm}

\textbf{重要的AI教育工具：}

\begin{table}[h]
\centering
\tablesize
\begin{tabular}{p{3.8cm}p{6.8cm}}
\toprule
\textbf{工具名稱} & \textbf{功能與特色} \\
\midrule
Affinity Learning & 免費平台,透過互動式情境模擬真實世界經驗 \\
& \url{https://affinitylearning.ca} \\
\midrule
IU Bloomington Simulation Scenario Builder Tool & 協助模擬教育工作者建立醫療保健模擬案例,免費提供 \\
& \url{https://chatgpt.com/g/g-mw2EfdWEj} \\
\midrule
AI Patient Actor App (Dartmouth) & 讓醫療學員練習臨床推理和溝通技巧,提供即時回饋 \\
& \url{https://geiselmed.dartmouth.edu/thesen/patient-actor-app/} \\
\bottomrule
\end{tabular}
\end{table}
\end{frame}

\begin{frame}{其他創新教育工具}
\frametitle{其他創新教育工具}

\textbf{Paper-based Simulation：}
\begin{itemize}
    \item 會議中討論了紙本模擬的應用
    \item 展示了不依賴高科技設備也能進行有效的臨床情境訓練
    \item 成本低、易於實施、適合資源有限的環境
\end{itemize}

\vspace{0.5cm}

\textbf{Standardized Patient的專業化：}
\begin{itemize}
    \item 芬蘭同學分享他們的經驗：
    \begin{itemize}
        \item 聘請專業演員擔任standardized patient
        \item 提供專業的表演課程訓練
        \item 提升模擬訓練的真實性和教育效果
    \end{itemize}
    \item 台大醫院的SP也有上專業的表演課程訓練,顯示這已成為國際趨勢
\end{itemize}

\vspace{0.3cm}

\begin{alertblock}{重要啟發}
科技工具與傳統方法的結合,創造更豐富的學習體驗。重視人文關懷和批判思考。
\end{alertblock}
\end{frame}

% 其他重要議題
\section{其他重要議題}

\begin{frame}{其他重要研究議題}
\frametitle{其他重要研究議題}

\textbf{Failure to Fail研究：}
\begin{itemize}
    \item \textbf{研究問題}：什麼時候我們會讓該被當掉的學生通過？
    \item 探討評估者為何難以給予不及格評分
    \item 影響評分決策的因素
    \item 如何建立公平且嚴謹的評估機制
\end{itemize}

\vspace{0.5cm}

\textbf{Visual Thinking Strategies (VTS)：}
\begin{itemize}
    \item 基於證據的教學方法,透過藝術分析和結構化探究
    \item 培養觀察能力、發展批判性思考
    \item 提升溝通技巧、促進團隊討論
    \item 參考文獻：\url{https://bmcmededuc.biomedcentral.com/articles/10.1186/s12909-025-06642-9}
\end{itemize}

\vspace{0.5cm}

\textbf{Simulation Scenario Development：}
\begin{itemize}
    \item 使用Simulation Scenario Evaluation Tool (SSET)評估情境品質
    \item 參考：\url{https://pmc.ncbi.nlm.nih.gov/articles/PMC8936988/}
    \item 確保情境設計符合學習目標,提供真實的臨床挑戰
\end{itemize}
\end{frame}

\begin{frame}{Discourse（話語/論述）}
\frametitle{Discourse（話語/論述）}

\textbf{多層面的意義：}

\begin{columns}
\begin{column}{0.5\textwidth}
\textbf{語言學/溝通學意義}
\begin{itemize}
    \item 超越單句的語言使用
    \item 包含完整的對話、文章或演說
    \item 有連貫性、結構性的語言表達
\end{itemize}

\vspace{0.3cm}

\textbf{社會學/哲學意義}
\begin{itemize}
    \item 特定社群或領域的話語體系和思維模式
    \item 反映權力關係、社會價值觀和知識結構
    \item 例：醫學discourse、法律discourse
\end{itemize}
\end{column}

\begin{column}{0.5\textwidth}
\textbf{傅柯(Foucault)的Discourse理論}
\begin{itemize}
    \item discourse不只是語言,更是權力和知識的載體
    \item 塑造我們如何理解世界和自我
    \item 例如：
    \begin{itemize}
        \item 醫學discourse如何定義「正常」與「病態」
        \item 教育discourse如何建構「好學生」標準
    \end{itemize}
\end{itemize}

\vspace{0.3cm}

\textbf{在醫學教育中的應用}
\begin{itemize}
    \item 專業discourse：醫學術語和表達方式
    \item 教學discourse：臨床教學的話語框架
\end{itemize}
\end{column}
\end{columns}

\vspace{0.3cm}

\begin{block}{簡單說}
discourse就是「特定領域的說話和思考方式」,它不只是語言表面,更深層影響我們如何理解和建構專業知識。
\end{block}
\end{frame}


% 會議心得
\section{會議心得與反思}

\begin{frame}{教育理念的多元性}
\frametitle{教育理念的多元性}

AMEE 2025展現了醫學教育理論與實務的豐富多樣性：

\begin{itemize}
    \item \textbf{從建構式配合到對話式專業主義}
    \begin{itemize}
        \item 強調學習者中心的教育設計
        \item 重視系統性思維和整體性規劃
    \end{itemize}
    
    \item \textbf{從武器化專業主義的批判到包容性教育的倡議}
    \begin{itemize}
        \item 彰顯社會正義導向
        \item 警惕專業標準可能帶來的排他性
    \end{itemize}
    
    \item \textbf{從傳統研究典範到多元研究方法}
    \begin{itemize}
        \item 鼓勵理論創新
        \item 整合量化與質性研究取向
    \end{itemize}
\end{itemize}

\vspace{0.5cm}

\begin{alertblock}{核心訊息}
醫學教育需要多元視角、批判思考和持續創新,避免單一標準的僵化思維。
\end{alertblock}
\end{frame}

\begin{frame}{科技與教育的融合}
\frametitle{科技與教育的融合}

AI工具的快速發展為醫學教育帶來新機會：

\begin{columns}
\begin{column}{0.5\textwidth}
\textbf{AI帶來的優勢：}
\begin{itemize}
    \item 模擬情境建構更加便捷
    \item 個人化學習成為可能
    \item 評估方式更加多元
    \item 即時回饋提升學習效果
\end{itemize}
\end{column}

\begin{column}{0.5\textwidth}
\textbf{需要注意的事項：}
\begin{itemize}
    \item 仍需重視人文關懷
    \item 培養批判思考能力
    \item 避免過度依賴科技
    \item 保持教育的溫度
\end{itemize}
\end{column}
\end{columns}

\vspace{0.5cm}

\textbf{實用工具推薦：}
\begin{itemize}
    \item Affinity Learning - 免費互動式模擬平台
    \item IU Bloomington Simulation Builder - 情境建構工具
    \item AI Patient Actor App - 臨床溝通練習
\end{itemize}

\vspace{0.3cm}

\begin{block}{平衡觀點}
善用科技工具提升教學效能,但不忘教育的人文本質和社會責任。
\end{block}
\end{frame}

\begin{frame}{評估公平性的重要性}
\frametitle{評估公平性的重要性}

從對比效應到利益衝突管理,會議強調：

\begin{enumerate}
    \item \textbf{評估過程中的偏誤識別與控制}
    \begin{itemize}
        \item 認識對比效應、光環效應等評估偏誤
        \item 建立標準化評分機制
        \item 進行評審員訓練和校準
    \end{itemize}
    
    \item \textbf{建立透明且公平的評估機制}
    \begin{itemize}
        \item 明確評估標準和程序
        \item 多元評估方法的運用
        \item 提供申訴和複核機制
    \end{itemize}
    
    \item \textbf{重視多元觀點和文化敏感度}
    \begin{itemize}
        \item 避免單一文化標準
        \item 考量不同背景學生的特殊需求
        \item 建立包容性的評估環境
    \end{itemize}
    
    \item \textbf{持續改進評估品質}
    \begin{itemize}
        \item 定期檢討評估機制
        \item 收集多方回饋
        \item 基於證據進行改善
    \end{itemize}
\end{enumerate}
\end{frame}

\begin{frame}{對台灣醫學教育的啟示}
\frametitle{對台灣醫學教育的啟示}

\begin{enumerate}
    \item \textbf{加強對話式專業主義教育}
    \begin{itemize}
        \item 從規定性標準轉向反思性對話
        \item 培養批判思考能力
        \item 建立包容性的學習環境
    \end{itemize}
    
    \item \textbf{重視評估公平性}
    \begin{itemize}
        \item 建立完善的評估品質管控機制
        \item 減少各種偏誤
        \item 確保利益衝突管理
    \end{itemize}
    
    \item \textbf{善用科技工具}
    \begin{itemize}
        \item 積極探索AI等新科技在教學與評估中的應用
        \item 發展本土化的教育工具
        \item 平衡科技與人文
    \end{itemize}
    
    \item \textbf{促進包容性教育}
    \begin{itemize}
        \item 警惕專業主義標準可能帶來的排他性
        \item 建立更包容的學習環境
        \item 尊重多元文化背景
    \end{itemize}
    
    \item \textbf{發展本土教育研究}
    \begin{itemize}
        \item 不僅引進國際經驗
        \item 更要發展符合本地脈絡的教育實踐
        \item 建立理論與實務的連結
    \end{itemize}
\end{enumerate}
\end{frame}

% 結論
\section{結論}

\begin{frame}{核心訊息}
\frametitle{AMEE 2025核心訊息}

\begin{block}{醫學教育需要}
\begin{itemize}
    \item \textbf{對話而非處方}：透過持續對話建構專業主義
    \item \textbf{公平而非形式}：重視評估的公平性與文化敏感度
    \item \textbf{包容而非排斥}：警惕專業標準的武器化傾向
    \item \textbf{創新而非僵化}：善用科技工具提升教學品質
    \item \textbf{反思而非接受}：培養批判性思考與終身學習能力
\end{itemize}
\end{block}

\vspace{0.5cm}

\textbf{未來展望：}
\begin{itemize}
    \item 持續改進醫學教育,培養具有專業能力、人文關懷與社會責任的優秀醫師
    \item 建立更公平、包容、有效的教育評估體系
    \item 善用科技創新,同時保持教育的溫度和深度
    \item 發展本土化的教育理論與實踐
\end{itemize}
\end{frame}

\begin{frame}{致謝}
\frametitle{致謝}

感謝本院支持參與AMEE 2025國際會議,此次學習經驗豐富且深具啟發性。

\vspace{0.5cm}

\textbf{特別感謝：}
\begin{itemize}
    \item AMEE大會主辦單位提供優質的學習平台
    \item 各Workshop講師的精彩分享與討論
    \item 國際同儕的交流與經驗分享
    \item 本院同仁對醫學教育的持續支持
\end{itemize}

\vspace{0.5cm}

期待將這些寶貴經驗轉化為實際行動,為提升本院醫學教育品質而努力。

\vspace{1cm}

\begin{center}
\Large
\textbf{謝謝聆聽！}

%\vspace{0.5cm}

\normalsize
謝慕揚, MD, PhD, FESC\\
台大醫院新竹分院教學部副主任\\
2025年
\end{center}
\end{frame}

\end{document}